\scanpage{150}
\mspage{24}
% archive sheet number: 149

Revenons à la majoration cardinale des
$X_\alpha(f)$, où $f:X\unsure{\text{---}|\longrightarrow}Y$ \unsure{décrits|décrit|reste} dans $\mathrm{Fl}^{\unsure{s}}M$, i.e.
$X,Y\in\mathrm{Filt}^s$. On suppose de plus $C_0\subset\deleted{\unsure{\mathrm{Filt}}}\mathrm{Fl}^rM$,
\deleted{\unsure{et}}
\[
\left\{\begin{aligned}
&C_0\subset\mathrm{Fl}^rM=\mathrm{Fl}\,M^r,\quad
(f\unsure{:|=}X\text{---}Y)\in\mathrm{Fl}^sM\quad\text{i.e. }X,Y\in\mathrm{Filt}^s\\
&\text{avec }s\geq r
\end{aligned}\right.
\tag*{(3.10)}
\]
\deleted{\unsure{Alors} pour $C_0$, \unsure{notons qu'aux.|notons qu'en} d'objets}
\[
\deleted{\operatorname{Card}\Delta(C_0,f)\leq\operatorname{Card}M^r_{/X}\times\operatorname{Card}M^r_{/Y}}
\]

Il y a une difficulté idiote, c'est que
\deleted{\unsure{de ce}} $M^r$ est équivalente à une petite catégorie,
mais n'est pas elle-même petite. Même si
$C_0$ ne contient qu'une seule classe d'isomorphie de flèches de $M^r$, son cardinal pourrait
être (pour $r$ fixé) arbitrairement élevé. Aussi
dans ce qui suit, je \inserted{dois} supposer \deleted{que le}
\inserted{la condition} \deleted{nb d'éléments, cardinal de \unsure{la}}
de ``\unsure{non pléthoricité} cardinale'' pour $C_0$

\[
\begin{aligned}
&\forall i\in C_0,\quad\deleted{\text{l'ens des}}\inserted{\text{désignant}} \text{ par }C_0(i)\text{ l'ens des}\\
&\qquad\deleted{\unsure{A}}i'\in C_0\text{ telles que }i\simeq i'\text{ dans }\mathrm{Fl}\,M,\text{ j'impose}\\
&\qquad\text{que }\operatorname{Card}C_0(i)\leq 2^{(r^{\pi_0})}
\end{aligned}
\tag*{(3.1\unsure{1|b})}
\]
(ce qui est \inserted{aussi} la \inserted{valeur de la} majoration pour le cardinal de
$\widetilde{\underline{\mathrm{Fl}}^rM}=\widetilde{\underline{\mathrm{Fl}}\,M^r}$)

Cela entraîne donc \deleted{que} l'on a
$\inserted{\operatorname{Card}C_0\leq}\deleted{\unsure{\operatorname{Card}}}(2^{(r^{\pi_0})})^2=\operatorname{Card}2^{r^{\pi_0}}$.
\[
\operatorname{Card}C_0\leq 2^{(r^{\pi_0})},
\tag*{(3.12)}
\]
qui est équivalente à (3.1\unsure{1|b}) (moyennant $C_0\subset\mathrm{Fl}^rM$,
postulée dans (3.10)).

Il s'ensuit \inserted{par (3.8)} que l'on a
\[
\operatorname{Card}\Delta(C_0,f)\leq
\underbrace{\operatorname{Card}(C_0)}_{\leq 2^{(r^{\pi_0})}}
\mathbin{\cdot}s^{(r^{\pi_0})}\mathbin{\cdot}s^{(r^{\pi_0})}\leq s^{(r^{\pi_0})}.
\]

D'autre part, les $X_D$ dont on prend la somme
amalgamée sous $X$, pour obtenir $X_1(f)$,
sont dans $M^s$, si on fait l'hypothèse suivante
